\section{Felhasználók és fiókkezelés}

Ez a fejezet részletesen bemutatja a felhasználók létrehozását, jellemző mezőit, jelszópolitika kezelését és a \texttt{/etc/passwd}, \texttt{/etc/shadow} fájlok felépítését.

\subsection{Felhasználó létrehozása és opciók}
\begin{lstlisting}[language=bash]
# Create a basic user with home directory and bash shell
useradd -m -s /bin/bash -G users user1
# Set a custom UID
useradd -u 1500 -m user2
\end{lstlisting}

Magyarázat:
\begin{itemize}
  \item \texttt{-m}: home létrehozása
  \item \texttt{-s /bin/bash}: beállítja a login shellt
  \item \texttt{-G users}: további csoporttagság
  \item \texttt{-u UID}: explicit UID számozás (ügyeljünk az ütközésekre)
\end{itemize}

\subsection{/etc/passwd és /etc/shadow}
\begin{verbatim}
# /etc/passwd line format
username:x:UID:GID:comment:home:shell
# /etc/shadow line format
username:passwordhash:lastchg:min:max:warn:inactive:expire:
\end{verbatim}

Gyakori műveletek:
\begin{lstlisting}[language=bash]
# Set password
passwd user1
# Lock password
passwd -l user1
# Unlock password
passwd -u user1
# Set password aging
chage -d 2026-01-01 -M 90 -W 7 user1
# Check aging settings
chage -l user1
\end{lstlisting}

Biztonsági megjegyzések:
\begin{itemize}
  \item A jelszóhasheket soha ne osszuk meg.
  \item Javasolt hosszú minimális és kötelező jelszópolitika, valamint többfaktoros autentikáció bevezetése.
\end{itemize}
